\documentclass[11pt]{article}
\usepackage[margin=0.9in]{geometry}
\usepackage{amsmath,amssymb,amsthm,booktabs}
\usepackage[T1]{fontenc}
\usepackage{lmodern,microtype}
\usepackage[colorlinks=true,linkcolor=blue,urlcolor=blue]{hyperref}
\newtheorem{theorem}{Theorem}[section]
\newtheorem{lemma}[theorem]{Lemma}
\newtheorem{proposition}[theorem]{Proposition}
\newtheorem{corollary}[theorem]{Corollary}
\theoremstyle{remark}
\newtheorem{remark}[theorem]{Remark}
\newcommand{\R}{\mathbb R}
\newcommand{\C}{\mathbb C}
\newcommand{\norm}[1]{\left\|#1\right\|}
\newcommand{\diag}{\operatorname{diag}}
\newcommand{\ran}{\operatorname{ran}}
\newcommand{\HH}{\mathcal H}
\title{Toward a four-dimensional Yang--Mills construction\\
\large Notebook 10: local curvature perturbations at a fixed constraint}
\author{Research working document}
\date{16 October 2026}
\begin{document}
\maketitle
\begin{abstract}
We identify the orthogonal projection onto feasible curvature
variations as two curl derivatives of the constrained flat inverse.
The full inverse estimates of Notebook 09 prove that this projection
is exponentially local in physical cell Hilbert norm, uniformly in
refinement and compatible periodic volume. This gives a convergent
weighted expansion for small local perturbations of the curvature
metric, with explicit factorial and source-path bounds at every order.
The averaging constraint is held fixed throughout. Identifying all
coefficients of an actual nonflat Yang--Mills background remains a
separate problem.
\end{abstract}

\noindent\textbf{Status.}
Research proof checked by a separate internal derivation and audit.
No external certification or novelty claim is made. The homogeneous
flat input operators are those matched to the pinned source account
in Notebooks 08 and 09. Fine spacing is $M^{-1}$ for odd $M$,
physical coarse spacing is one, and dimension is fixed with $d=4$
the intended application. All estimates below use the same physical
normalizations as those notebooks.
% Fixed-Q curvature perturbations; no standalone preamble.
% Requires amsmath/amssymb/amsthm, the existing theorem environments,
% and the existing \norm macro. Labels use the unique prefix curv-.
\section{Local curvature perturbations at a fixed averaging constraint}
\label{curv-section}

This section concerns a flat linear problem with its averaging map held
fixed. It supplies a full weighted analytic neighborhood for perturbing
the curvature quadratic form. It does not identify a non-Abelian
background variation, which can also change the averaging constraint,
covariant derivatives, gauge projection, lower-order terms, and current.

Let $D$ be the fine scalar gradient, $d$ the fine curl, and
$\mathcal L=D^*D$. Fix $Q=Q_M^{\rm st}$ and its scalar counterpart
$Q_0$, so that $QD=d_cQ_0$. All fine Hilbert norms use physical fine
volume, and the coarse cell norm uses counting measure. Retain the
matched scalar source projection $R$ and scalar inverse $G'$, with
\begin{equation}\label{curv-scalar-identities}
 Q_0G'R=0,\qquad \mathcal L G'R=R.
\end{equation}
For a fixed positive penalty $a$, let
\begin{equation}\label{curv-fine-data}
 \mathcal A_a=d^*d+DRD^*+Q^*aQ,\qquad G_0=\mathcal A_a^{-1},
 \qquad q_{\rm aux}=(QG_0Q^*)^{-1}=q_M+aI.
\end{equation}
Here $q_M$ is the physical Maxwell minimum and $q_{\rm aux}$ is the
distinct penalized coarse precision. The matched unit-scale source
has $a=1$. Let
\begin{equation}\label{curv-extension}
 H=G_0Q^*q_{\rm aux}=H_{\rm src},\qquad F=dH.
\end{equation}
Thus $QH=I$, $RD^*H=0$, and $q_M=F^*F$.
The full inverse theorem supplies uniform exponential cell $L^2$
locality for $G_0$ and all its fine derivative sandwiches of total
order at most two. The minimizing extension theorem supplies the same
cell locality for $F$. These inputs hold uniformly in odd $M$ and on
every compatible periodic torus, as well as in homogeneous infinite
volume. No assertion for general Dirichlet boundaries is used here.

For cell-operator kernels, use the weighted algebra norm
\begin{equation}\label{curv-schur-norm}
 \norm{T}_\kappa=
 \max\left\{\sup_x\sum_y e^{\kappa\operatorname{dist}(x,y)}
                  \norm{T(x,y)},\quad
             \sup_y\sum_x e^{\kappa\operatorname{dist}(x,y)}
                  \norm{T(x,y)}\right\}.
\end{equation}
The distance is $\ell^1$ cell distance, or its periodic version.
The internal norms in \eqref{curv-schur-norm} are operator norms
between the physical-volume cell Hilbert spaces; this convention
also applies to maps between coarse bond cells and fine curvature
cells. The norm is submultiplicative and unchanged by adjoints.

\begin{theorem}\label{curv-theorem}
Let $\mathcal K=\overline{d\ker Q}$ in the curvature Hilbert space,
and let $P_{\mathcal K}$ be its orthogonal projection. Then
\begin{equation}\label{curv-projection}
 \Gamma=G_0-Hq_{\rm aux}^{-1}H^*,\qquad
 P_{\mathcal K}=d\Gamma d^*
   =dG_0d^*-Fq_{\rm aux}^{-1}F^*.
\end{equation}
For a sufficiently small fixed $\kappa>0$, the constants
\begin{equation}\label{curv-local-constants}
 f=\norm{F}_\kappa,\qquad
 p=\max\left\{1,\norm{dG_0d^*}_\kappa
                  +f^2\norm{QG_0Q^*}_\kappa\right\}
\end{equation}
are bounded independently of refinement and compatible periodic
volume, and $\norm{P_{\mathcal K}}_\kappa\le p$.

For any complex curvature operator $V$ satisfying
$p\norm{V}_\kappa<1$, the expression
\begin{equation}\label{curv-holomorphic-action}
 q(V)=q_M+F^*V(I+P_{\mathcal K}V)^{-1}F
\end{equation}
is a holomorphic weighted-cell-operator map. On the real selfadjoint
slice, it equals the constrained positive-form minimum
\begin{equation}\label{curv-real-minimum}
 \langle B,q(V)B\rangle
   =\inf_{QA=B}\langle dA,(I+V)dA\rangle.
\end{equation}
It obeys
\begin{align}
 \norm{q(V)-q_M}_\kappa
   &\le\frac{f^2\norm{V}_\kappa}{1-p\norm{V}_\kappa},
       \label{curv-first-bound}\\
 \norm{q(V)-q_M-F^*VF}_\kappa
   &\le\frac{f^2p\norm{V}_\kappa^2}{1-p\norm{V}_\kappa}.
       \label{curv-remainder-bound}
\end{align}
In particular, for disjoint cell-supported operators
$V=\bigoplus_x V_x$, the condition is the volume-independent bound
$p\sup_x\norm{V_x}<1$.
\end{theorem}

\begin{proof}
The identities $QH=I$ and $G_0\mathcal A_a=I$ give
\begin{equation}\label{curv-constrained-inverse}
 Q\Gamma=\Gamma Q^*=0,\qquad
 \Gamma\mathcal A_a=I-HQ.
\end{equation}
Consequently $\Gamma$ is the inverse of the restricted positive
form on $\ker Q$, extended by zero on its orthogonal complement
$\operatorname{ran}Q^*$. The latter range is closed because $Q$
has the bounded right inverse $H$.

For any curvature vector $u$, put $\zeta=\Gamma d^*u$ and
$r=RD^*\zeta$. The constrained Euler equation is
\begin{equation}\label{curv-euler}
 \langle d\zeta-u,d\xi\rangle
    +\langle RD^*\zeta,RD^*\xi\rangle=0,
       \qquad \xi\in\ker Q.
\end{equation}
Set $\lambda=G'r=G'Rr$. Equation
\eqref{curv-scalar-identities} gives $Q_0\lambda=0$ and
$\mathcal L\lambda=r$. Thus $\xi=D\lambda$ is an admissible
test: $Q\xi=d_cQ_0\lambda=0$, $d\xi=0$, and
$RD^*\xi=Rr=r$. Equation \eqref{curv-euler} yields
$\norm{r}^2=0$. It then says that $d\zeta\in\mathcal K$ and
$u-d\zeta\perp\mathcal K$, proving
$d\Gamma d^*u=P_{\mathcal K}u$. At each fixed refinement all
difference operators are bounded, so the same test is valid in
infinite volume. No refinement-uniform norm for an individual
fine derivative is needed.

Equation \eqref{curv-projection} now follows by expanding $\Gamma$.
Its locality estimate follows from the two-derivative inverse bound,
the locality of $F$, and
$q_{\rm aux}^{-1}=QG_0Q^*$. In particular, no inverse of the
massless precision $q_M$ is taken. Submultiplicativity gives
\eqref{curv-local-constants}.

Flat minimality implies $P_{\mathcal K}F=0$. For a real
selfadjoint $V$ with $\norm{V}<1$, every feasible curvature lies
in the affine space $FB+d\ker Q$, whose closure is
$FB+\mathcal K$. Passing to this closure does not change the
infimum of the bounded strictly positive form $I+V$.
The minimizing correction $k\in\mathcal K$ satisfies
\[
 (I_{\mathcal K}+P_{\mathcal K}VP_{\mathcal K})k
       =-P_{\mathcal K}VFB.
\]
Substitution gives
\begin{equation}\label{curv-schur-action}
 q(V)=q_M+F^*VF-
 F^*VP_{\mathcal K}(I_{\mathcal K}
       +P_{\mathcal K}VP_{\mathcal K})^{-1}
       P_{\mathcal K}VF.
\end{equation}
Expanding the inverse and using $P_{\mathcal K}^2=P_{\mathcal K}$
identifies this with \eqref{curv-holomorphic-action}. Both $V$
factors in \eqref{curv-schur-action} remain $V$ in the complex
continuation; replacing one by $V^*$ would destroy holomorphy.
The Neumann series converges in the weighted algebra for
$p\norm{V}_\kappa<1$ and proves
\eqref{curv-first-bound}--\eqref{curv-remainder-bound}.
This condition implies $\norm{V}<1$, since $p\ge1$, on the real
selfadjoint slice. Finally a cell-diagonal direct sum has
$\norm{V}_\kappa=\sup_x\norm{V_x}$.
\end{proof}

\begin{corollary}[All-order anchored curvature responses]
\label{curv-all-orders}
For $n\ge1$,
\begin{equation}\label{curv-derivative-formula}
 D^nq(0)[V_1,\ldots,V_n]
   =(-1)^{n-1}\sum_{\sigma\in S_n}
       F^*V_{\sigma(1)}P_{\mathcal K}V_{\sigma(2)}
       \cdots P_{\mathcal K}V_{\sigma(n)}F,
\end{equation}
with $n-1$ projection factors. Hence
\begin{equation}\label{curv-derivative-norm}
 \norm{D^nq(0)[V_1,\ldots,V_n]}_\kappa
       \le n!f^2p^{n-1}\prod_{j=1}^n\norm{V_j}_\kappa.
\end{equation}
If $V_j$ is supported wholly in curvature cell $x_j$, then the
block from coarse cell $b$ to coarse cell $c$ obeys
\begin{align}
 &\norm{D^nq(0)[V_1,\ldots,V_n](b,c)}\notag\\
 &\quad\le f^2p^{n-1}\prod_{j=1}^n\norm{V_j}
       \sum_{\sigma\in S_n}
       \exp\!\left[-\kappa\left(
          \operatorname{dist}(b,x_{\sigma(1)})
          +\sum_{i=1}^{n-1}\operatorname{dist}
                   (x_{\sigma(i)},x_{\sigma(i+1)})
          +\operatorname{dist}(x_{\sigma(n)},c)
                       \right)\right].
                       \label{curv-anchored-bound}
\end{align}
Repeated source cells are allowed, and the same constants apply
on compatible tori with periodic distance.
\end{corollary}

\begin{proof}
Differentiate the Neumann series in
\eqref{curv-holomorphic-action}. Its homogeneous degree-$n$ term
has $n-1$ projections and the sign displayed in
\eqref{curv-derivative-formula}; polarization sums the $n!$
orders. Equation \eqref{curv-derivative-norm} follows by the algebra
norm. Each block of an operator of weighted norm $c$ is bounded
by $ce^{-\kappa\operatorname{dist}}$. Applying this to the two
$F$ end legs and the $n-1$ projection links gives
\eqref{curv-anchored-bound}. The cell-supported insertions fix
every intermediate cell, so there is no spatial convolution
or unnormalized sum over internal fine coordinates.
\end{proof}

The cell-diagonal norm identity above uses disjoint cell subspaces.
Overlapping supports require their actual weighted operator norm
or an explicit overlap estimate. The theorem controls only the
curvature metric at fixed $Q$; it does not supply the other terms
generated by a nonflat Yang--Mills background variation.

\section{Use in the construction program}

A nonflat background changes more than the coefficient of the
curvature quadratic form. The present theorem isolates a component
whose complete response is now controlled. To apply it to an actual
source family, one must identify a curvature coefficient $V$ in the
stated weighted ball and handle the accompanying changes to the
covariant derivatives, scalar gauge, corrected average and lower-order
terms. A finite-regulator derivative formula alone does not verify
that common domain. The auxiliary-current domain of Notebook 09 is a
separate proved perturbation direction.

The conclusion concerns the local quadratic reference theory. It
does not produce the interacting ultraviolet measure, restore all
continuum symmetries, or establish nontriviality or a physical mass gap.
Those obligations remain as recorded in Notebook 01.

\begin{thebibliography}{9}
\bibitem{repo}
\href{https://github.com/ethanwsussman/Yang--Mills-LLM-experiment/tree/e392ef45013a581664bab76b90867170b30fe10e}
{Pinned Yang--Mills-LLM-experiment account}, commit
\texttt{e392ef45013a581664bab76b90867170b30fe10e}.
\bibitem{previous}
Companion research notebooks: Notebook 07,
\emph{Flat covariance locality and its free refinement limit};
Notebook 08, \emph{A local flat minimizing extension}; and Notebook 09,
\emph{The flat fine inverse and a current neighborhood}.
Their internal proofs supply the homogeneous operator estimates
used here; they do not certify the complete deposited manuscript.
\end{thebibliography}
\end{document}